\phantomsection  % so hyperref creates bookmarks
\addcontentsline{toc}{chapter}{Información sobre Licencia}
%\label{label_fdl}

 \begin{center}

       Información sobre Licencia


\end{center}

El proyecto combina dos tipos de obra con necesidades de licencia distintas: la memoria y el material gráfico que la acompaña forman la documentación del TFG, mientras que los motores de generación, el servidor y la aplicación móvil son el software desarrollado. Se adoptan dos licencias diferentes para que cada parte se redistribuya en los términos que mejor se ajustan a su naturaleza y a las dependencias que incorpora.

La memoria y el resto de documentación asociada al TFG se distribuyen bajo la licencia Creative Commons Reconocimiento-NoComercial-CompartirIgual 4.0 Internacional (CC BY-NC-SA 4.0)~\cite{ccbyncsa4}. Esta licencia permite copiar, redistribuir y adaptar el material siempre que se atribuya la autoría, no se haga un uso comercial y las obras derivadas se publiquen bajo la misma licencia. La elección responde a dos razones. La memoria incorpora pictogramas del catálogo ARASAAC~\cite{arasaac}, que se distribuyen bajo CC BY-NC-SA 4.0~\cite{ccbyncsa4}; la cláusula de CompartirIgual obliga a redistribuir cualquier obra derivada de estos pictogramas bajo la misma licencia, de modo que aplicar CC BY-NC-SA 4.0 al conjunto documental es la opción coherente con el material embebido. La cláusula NoComercial, además, encaja con el carácter académico del trabajo y con el espíritu del catálogo ARASAAC, cuyo uso está pensado para fines educativos y terapéuticos más que para explotación comercial.

El software desarrollado en el proyecto (los seis motores de generación, el \textit{pipeline} de evaluación, el servidor \texttt{FastAPI} y la aplicación móvil Dixi) se distribuye bajo la licencia MIT~\cite{mitlicense}. Es una licencia permisiva que autoriza el uso, la modificación y la redistribución del código, con o sin fines comerciales, siempre que se conserven el aviso de copyright y el texto de la licencia. Esta elección persigue compatibilidad y reutilización en igual medida. La MIT es compatible con prácticamente la totalidad de las bibliotecas de terceros que el proyecto utiliza (\texttt{React Native}, \texttt{Expo}, \texttt{FastAPI}, \texttt{Lark}, \texttt{spaCy}, \texttt{Transformers} de \texttt{Hugging Face}), lo que evita conflictos de licencia en despliegues derivados. Una licencia permisiva facilita que otros estudiantes, investigadores o desarrolladores del ámbito de CAA partan del código de Dixi para construir sistemas propios sin asumir las obligaciones de una copyleft. Dado que el valor del proyecto reside más en las decisiones metodológicas y en los datos consolidados que en un código reutilizable cerrado, la prioridad es eliminar barreras a su reutilización.

Cada licencia se distribuye por separado. El texto íntegro de CC BY-NC-SA 4.0 está disponible en la página oficial de Creative Commons~\cite{ccbyncsa4}, y el texto íntegro de la licencia MIT, en la Open Source Initiative~\cite{mitlicense}. Cada repositorio de código incluye el fichero de licencia correspondiente junto al resto del material, y la memoria reserva esta declaración para documentar la elección de forma visible al lector.

Quedan respetados los derechos de terceros en todos los casos. Los pictogramas ARASAAC mantienen su licencia original CC BY-NC-SA 4.0, con atribución al Gobierno de Aragón como titular del catálogo~\cite{arasaac}, y las dependencias de software conservan las licencias de sus autores originales, documentadas en los ficheros de licencia de cada paquete. Las referencias bibliográficas citadas a lo largo de la memoria se acogen al derecho de cita previsto en la legislación vigente sobre propiedad intelectual. Ningún contenido de este TFG se apropia de material protegido sin la debida atribución o autorización.
